\subsection{Channel-doping sweep: completed experiment}
With the student-specific parameter $\delta=2$, the Chapter 4
channel-doping values are $N_{\mathrm{CH}}=3\times10^{15}$,
$3\times10^{16}$, $3\times10^{17}$, and
$3\times10^{18}$ cm$^{-3}$. The gate length and upper-silicon
thickness were fixed at $L_G=50$ nm and $T_{\mathrm{SI}}=30$ nm.
The other SProcess settings were held at the supplied nominal values,
including $E_{\mathrm{EXT}}=0.5$ keV, $E_{\mathrm{SD}}=15$ keV
and $T_{\mathrm{ANNEAL}}=950\,^\circ$C. Four SProcess runs
and eight SDevice $I_D$--$V_G$ runs completed.
Each device was swept at $V_{DS}=0.1$ V and 1.2 V.
The associated process/device node groups were
$n1/n2,n14$; $n4/n5,n15$; $n6/n7,n16$; and
$n8/n9,n17$, in order of increasing doping.

\paragraph{Extraction method.}
The eight Sentaurus Visual exported current traces were reduced to
35 unique gate-bias samples per curve (final value retained if the
drain-initialization ramp reused $V_{GS}=-0.5$ V).
The cleaned, matched $I_D$--$V_{GS}$ dataset is stored in
\texttt{data/part4\_idvg\_all.csv}.
The operating-point proxies are
$I_{\mathrm{ON}}=|I_D|(V_{GS}=1.2\mathrm{\,V},V_{DS}=1.2\mathrm{\,V})$
and
$I_{\mathrm{OFF}}=|I_D|(V_{GS}=0,V_{DS}=1.2\mathrm{\,V})$.
The current magnitudes retain the native exported simulator
units without applying another width conversion. The supplied
baseline SDevice script specifies \texttt{AreaFactor=1.0},
and the modified Workbench decks are assumed to retain this setting. The constant-current $V_{\mathrm{TH,CC}}$ reference
is $|I_D|=10^{-7}$ in those units, with linear interpolation
in $\log_{10}|I_D|$ between simulated samples. No threshold
is extrapolated beyond the $-0.5$ to $1.2$ V gate sweep.
The reported apparent SS uses the common $-0.20$ to
$+0.05$ V, $V_{DS}=0.1$ V fit window for
$\log_{10}|I_D|$; as in Parts 2 and 3, a strong regression fit
does not prove true weak inversion in this window.
DIBL uses matched thresholds
$\mathrm{DIBL}_{\mathrm{CC}} =
1000[V_{\mathrm{TH,CC}}(0.1)-V_{\mathrm{TH,CC}}(1.2)]/1.1$
in mV/V, and is N/A where either threshold cannot be extracted.

\begin{table}[H]\centering\small
\caption{Part 4 metrics from the original Sentaurus CSV exports.
N/A means the selected constant-current threshold is outside the
simulated gate-voltage range, not that a simulation failed.
$I_{\mathrm{ON}}$ and $I_{\mathrm{OFF}}$ use simulator-exported units.}
\resizebox{\linewidth}{!}{%
\begin{tabular}{@{}rrrrrrrr@{}}\toprule
$N_{\mathrm{CH}}$ (cm$^{-3}$) &
$V_{\mathrm{TH,0.1}}$ (V) &
$V_{\mathrm{TH,1.2}}$ (V) &
SS (mV/dec) & $I_{\mathrm{ON}}$ &
$I_{\mathrm{OFF}}$ & $I_{\mathrm{ON}}/I_{\mathrm{OFF}}$ &
DIBL (mV/V)\\\midrule
$3\times10^{15}$ & N/A & N/A & 351.80 & $1.44103\times10^{-3}$ & $1.58466\times10^{-4}$ & 9.09 & N/A\\
$3\times10^{16}$ & N/A & N/A & 334.20 & $1.42853\times10^{-3}$ & $1.46857\times10^{-4}$ & 9.73 & N/A\\
$3\times10^{17}$ & $-0.18582$ & N/A & 238.64 & $1.28887\times10^{-3}$ & $7.67334\times10^{-5}$ & 16.80 & N/A\\
$3\times10^{18}$ & $0.48383$ & $0.38090$ & 81.35 & $5.49280\times10^{-4}$ & $1.72411\times10^{-12}$ & $3.19\times10^8$ & 93.57\\
\bottomrule
\end{tabular}}
\end{table}

\paragraph{Threshold extraction and all-case switching-voltage comparison.}
The original $10^{-7}$-current constant-current definition remains the
reported threshold-voltage method. It is genuinely not extractable
for several low-doping/bias cases within the simulated gate range:
those values remain N/A in the table. To show a \emph{complete}
four-doping bias-dependent voltage comparison \emph{without extrapolation},
we additionally evaluate the gate voltage $V_{\mathrm{SW}}$ at
$|I_D|=2\times10^{-4}$ in simulator-exported units. All eight
measured curves cross this reference within the available sweep.
The crossing uses the same linear-in-$\log|I_D|$ interpolation
between adjacent sampled $V_{GS}$ points. This is a
\emph{high-current switching-voltage marker}, not a replacement
for physical threshold voltage or a subthreshold criterion.
The complete four-case marker data are archived in
\texttt{data/part4\_common\_current\_switching.csv}.
\begin{figure}[!htbp]\centering
\begin{tikzpicture}
\begin{semilogxaxis}[width=.89\linewidth,height=6.2cm,
xlabel={$N_{\mathrm{CH}}$ (cm$^{-3}$)},
ylabel={Gate-voltage crossing (V)},
xmin=1.5e15,xmax=5e18,ymin=-.35,ymax=1.30,
xtick={3e15,3e16,3e17,3e18},
xticklabels={$3\times10^{15}$,$3\times10^{16}$,$3\times10^{17}$,$3\times10^{18}$},
grid=both,legend style={font=\scriptsize,at={(0.02,.98)},anchor=north west}]
\addplot+[mark=*,thick]
table[col sep=comma,x=NCH_cm3,y=Vshift_0p1_V]
{data/part4_common_current_switching.csv};
\addlegendentry{$V_{\mathrm{SW}}$, $V_{DS}=0.1$ V}
\addplot+[mark=square*,thick,dashed]
table[col sep=comma,x=NCH_cm3,y=Vshift_1p2_V]
{data/part4_common_current_switching.csv};
\addlegendentry{$V_{\mathrm{SW}}$, $V_{DS}=1.2$ V}
\addplot+[only marks,mark=triangle*,mark size=2.4pt]
coordinates {(3e17,-.18581833)(3e18,.48382523)};
\addlegendentry{Original $V_{\mathrm{TH,CC}}$, 0.1 V}
\addplot+[only marks,mark=diamond*,mark size=2.4pt]
coordinates {(3e18,.38089576)};
\addlegendentry{Original $V_{\mathrm{TH,CC}}$, 1.2 V}
\end{semilogxaxis}
\end{tikzpicture}
\caption{Complete measured $V_{\mathrm{SW}}$ trends (curves)
at a common high-current crossing of $2\times10^{-4}$ exported
units, with the separately measured $10^{-7}$-criterion
threshold voltages (isolated symbols). The markers labelled
$V_{\mathrm{SW}}$ are \emph{not} threshold voltages; missing
$V_{\mathrm{TH,CC}}$ values are intentionally not extrapolated.}
\end{figure}

\begin{figure}[!htbp]\centering
\begin{tikzpicture}
\begin{semilogxaxis}[width=.90\linewidth,height=6.2cm,
xlabel={$N_{\mathrm{CH}}$ (cm$^{-3}$)},
ylabel={Apparent SS (mV/dec)},xmin=1.5e15,xmax=5e18,
xtick={3e15,3e16,3e17,3e18},
xticklabels={$3\times10^{15}$,$3\times10^{16}$,$3\times10^{17}$,$3\times10^{18}$},
grid=both]
\addplot+[mark=*,thick]
table[col sep=comma,x=NCH_cm3,y=SS_mV_dec]
{data/part4_doping_metrics.csv};
\end{semilogxaxis}
\end{tikzpicture}
\caption{Apparent subthreshold swing extracted over the same
$V_{GS}$ regression window at $V_{DS}=0.1$ V.}
\end{figure}

\begin{figure}[H]\centering
\begin{tikzpicture}
\begin{semilogxaxis}[width=.90\linewidth,height=6.2cm,
xlabel={$N_{\mathrm{CH}}$ (cm$^{-3}$)},
ylabel={$I_{\mathrm{ON}}$ (exported current units)},
xmin=1.5e15,xmax=5e18,
xtick={3e15,3e16,3e17,3e18},
xticklabels={$3\times10^{15}$,$3\times10^{16}$,$3\times10^{17}$,$3\times10^{18}$},
grid=both]
\addplot+[mark=*,thick]
table[col sep=comma,x=NCH_cm3,y=ION_exported]
{data/part4_doping_metrics.csv};
\end{semilogxaxis}
\end{tikzpicture}
\caption{ON-state current at $V_{GS}=V_{DS}=1.2$ V.}
\end{figure}

\begin{figure}[H]\centering
\begin{tikzpicture}
\begin{loglogaxis}[width=.90\linewidth,height=6.2cm,
xlabel={$N_{\mathrm{CH}}$ (cm$^{-3}$)},
ylabel={$I_{\mathrm{OFF}}$ (exported current units)},
xmin=1.5e15,xmax=5e18,ymin=1e-13,ymax=1e-3,
xtick={3e15,3e16,3e17,3e18},
xticklabels={$3\times10^{15}$,$3\times10^{16}$,$3\times10^{17}$,$3\times10^{18}$},
grid=both]
\addplot+[mark=*,thick]
table[col sep=comma,x=NCH_cm3,y=IOFF_exported]
{data/part4_doping_metrics.csv};
\end{loglogaxis}
\end{tikzpicture}
\caption{OFF-state current at $V_{GS}=0$, $V_{DS}=1.2$ V.
Both concentration and leakage are plotted logarithmically
to show the large dynamic range.}
\end{figure}

\begin{figure}[!htbp]\centering
\begin{tikzpicture}
\begin{semilogxaxis}[width=.89\linewidth,height=6.2cm,
xlabel={$N_{\mathrm{CH}}$ (cm$^{-3}$)},
ylabel={Drain-induced gate-voltage shift (mV/V)},
xmin=1.5e15,xmax=5e18,ymin=0,ymax=410,
xtick={3e15,3e16,3e17,3e18},
xticklabels={$3\times10^{15}$,$3\times10^{16}$,$3\times10^{17}$,$3\times10^{18}$},
grid=both,legend style={font=\scriptsize,at={(0.02,.98)},anchor=north west}]
\addplot+[mark=square*,thick,dashed]
table[col sep=comma,x=NCH_cm3,y=DrainShift_mV_per_V]
{data/part4_common_current_switching.csv};
\addlegendentry{$V_{\mathrm{SW}}$ shift (not DIBL)}
\addplot+[only marks,mark=*,mark size=3pt]
coordinates {(3e18,93.57225)};
\addlegendentry{Extractable strict-current DIBL}
\end{semilogxaxis}
\end{tikzpicture}
\caption{A complete \emph{high-current gate-voltage shift}
comparison (dashed curve) is displayed separately from
the physically distinct $10^{-7}$ reference-current DIBL
measurement (single filled circle, 93.57 mV/V).
The high-current values must \emph{not} be interpreted
as subthreshold DIBL; other strict-current DIBL
values are outside the measured sweep.}
\end{figure}

\paragraph{Required low-/high-doping transfer-curve overlay.}
At the common high drain bias $V_{DS}=1.2$ V, the two extreme
channel dopings are plotted with the same logarithmic ordinate.
These curves come directly from the cleaned exported transfer
traces; no fitted or extrapolated values are included.
\begin{figure}[H]\centering
\begin{tikzpicture}
\begin{semilogyaxis}[width=.93\linewidth,height=7.0cm,
xlabel={$V_{GS}$ (V)},ylabel={$|I_D|$ (exported current units)},
xmin=-0.5,xmax=1.2,ymin=1e-17,ymax=1e-2,
grid=both,legend pos=north west]
\addplot+[no marks,thick]
table[col sep=comma,x=VGS_V,y=Id_3e15_1p2]
{data/part4_idvg_all.csv};
\addlegendentry{$N_{\mathrm{CH}}=3\times10^{15}$ cm$^{-3}$}
\addplot+[no marks,thick]
table[col sep=comma,x=VGS_V,y=Id_3e18_1p2]
{data/part4_idvg_all.csv};
\addlegendentry{$N_{\mathrm{CH}}=3\times10^{18}$ cm$^{-3}$}
\end{semilogyaxis}
\end{tikzpicture}
\caption{Measured semilogarithmic transfer curves at
$V_{DS}=1.2$ V, comparing the extreme channel dopings.
The greater high-doping OFF suppression is accompanied by
a lower ON current.}
\end{figure}

\paragraph{Required channel-centre vertical electron density.}
Sentaurus Visual cutlines were exported from the two completed
high-drain SDevice results: node \texttt{n14\_des.tdr}
($3\times10^{15}$ cm$^{-3}$) and \texttt{n17\_des.tdr}
($3\times10^{18}$ cm$^{-3}$). Both correspond to the final
transfer-sweep operating point $V_{GS}=V_{DS}=1.2$ V.
The cutline used the lateral channel centre, $y=0.0$,
and sampled the vertical depth through the
30-nm top-silicon film. For the exported Sentaurus Visual
coordinate labelled \texttt{Y}, the top-silicon interval
runs approximately $-0.180$ to $-0.150\,\mu$m.
Depth is reported as
$d=1000(Y+0.180)\,$nm.
Only silicon records with nonzero acceptor concentration
were retained; the identical-coordinate oxide/interface
records with zero \texttt{eDensity} were excluded.
Both profiles retain 37 matched positive-concentration
silicon data points in
\texttt{data/part4\_electron\_vertical.csv}.

\begin{figure}[H]\centering
\begin{tikzpicture}
\begin{semilogyaxis}[width=.93\linewidth,height=7cm,
xlabel={Depth below top-Si gate interface (nm)},
ylabel={Electron density $n$ (cm$^{-3}$)},
xmin=0,xmax=30,ymin=1e6,ymax=1e21,
grid=both,legend pos=south west]
\addplot+[no marks,very thick]
table[col sep=comma,x=Depth_nm,y=eDensity_NCH3e15_cm3]
{data/part4_electron_vertical.csv};
\addlegendentry{$N_{\mathrm{CH}}=3\times10^{15}$ cm$^{-3}$}
\addplot+[no marks,very thick]
table[col sep=comma,x=Depth_nm,y=eDensity_NCH3e18_cm3]
{data/part4_electron_vertical.csv};
\addlegendentry{$N_{\mathrm{CH}}=3\times10^{18}$ cm$^{-3}$}
\end{semilogyaxis}
\end{tikzpicture}
\caption{Measured channel-centre vertical electron-density profiles
for the extreme doping cases at $V_{GS}=V_{DS}=1.2$ V.
The profile spans the upper-silicon film from the gate side
to the buried-oxide interface; oxide data were removed.}
\end{figure}

The electron density remains high at the gate-side silicon interface
under $V_{GS}=1.2$ V for both devices:
$9.77\times10^{19}$ and $9.48\times10^{19}$ cm$^{-3}$.
However, at $d\approx15.04$ nm the
low-doping result is $6.64\times10^{17}$ cm$^{-3}$,
whereas the high-doping result is
$8.41\times10^8$ cm$^{-3}$. The higher-doped silicon
therefore has strong carrier suppression over much of the film,
even under the stated positive gate bias.
The associated SDevice \texttt{eMobility} cutline
at the same depth decreases from approximately
$54.58$ to $37.95$ in simulator-exported mobility units.
This supports the expected mobility penalty in this
particular biased device, alongside the larger threshold
and body-charge effects. These profiles are for the stated
\emph{ON} gate bias; they must not be misrepresented as
measured OFF-state electron-density profiles.

\paragraph{Question: appreciable leakage reduction before severe drive-current loss.}
For the discrete sweep, define \emph{appreciable} as at least
a 50\% reduction in $I_{\mathrm{OFF}}$ relative to the
$3\times10^{15}$ cm$^{-3}$ case,
and \emph{severe} as a reduction of at least 30\%
in $I_{\mathrm{ON}}$ against the same reference,
with both values measured at the stated bias points.
At $3\times10^{17}$ cm$^{-3}$, OFF current falls from
$1.58466\times10^{-4}$ to $7.67334\times10^{-5}$
(51.58\% reduction), while ON current falls from
$1.44103\times10^{-3}$ to $1.28887\times10^{-3}$
(10.56\% penalty). This meets the improvement
criterion without a severe drive penalty.
At $3\times10^{18}$ cm$^{-3}$, ON current drops by
61.88\%, exceeding the severe-loss threshold despite
an extremely large OFF-leakage reduction.
Among the four \emph{tested} concentrations, only
$3\times10^{17}$ cm$^{-3}$ meets the defined leakage
improvement and drive-current retention criteria relative
to the lowest-doping reference. The intermediate
$3\times10^{16}$ cm$^{-3}$ case does not achieve a 50\%
leakage reduction, and the next tested concentration,
$3\times10^{18}$ cm$^{-3}$, incurs severe ON-current loss.
Thus the useful improvement is demonstrated at
$3\times10^{17}$ cm$^{-3}$; the boundaries of any continuous
acceptable doping interval cannot be determined from
only four simulated points.
